\documentclass[11pt,aspectratio=169]{beamer}

\usetheme{Madrid}
\usecolortheme{beaver}
\setbeamertemplate{navigation symbols}{}

\usepackage[utf8]{inputenc}
\usepackage[T1]{fontenc}
\usepackage{amsmath,amssymb}
\usepackage{booktabs}
\usepackage{listings}
\usepackage{xcolor}
\usepackage{tikz}
\usepackage{tcolorbox}
\tcbuselibrary{skins,breakable}

\definecolor{codebg}{RGB}{245,245,245}
\definecolor{codekw}{RGB}{0,0,180}
\lstdefinestyle{latexcode}{
  backgroundcolor=\color{codebg},
  basicstyle=\ttfamily\scriptsize,
  keywordstyle=\color{codekw}\bfseries,
  commentstyle=\color{gray},
  frame=single,
  rulecolor=\color{gray!40},
  breaklines=true,
  showstringspaces=false,
  xleftmargin=4pt,
  xrightmargin=4pt,
}
\lstset{style=latexcode}

% ---------- Output box: mimics the compiled PDF result ----------
\newtcolorbox{outputbox}{
  colback=white, colframe=black!65, boxrule=0.5pt, arc=1pt,
  left=8pt, right=8pt, top=5pt, bottom=4pt,
  title=\textsc{Compiled Output}, fonttitle=\bfseries\scriptsize,
  coltitle=white, colbacktitle=black!65
}

% ---------- Explanation lead-in ----------
\newcommand{\whybox}[1]{%
  \vspace{2pt}
  {\scriptsize\raggedright\textcolor{red!55!black}{\textbf{Why \& What:}}~#1\par}
}

\newtcolorbox{exercisebox}{
  colback=orange!8, colframe=orange!70!black, boxrule=0.6pt,
  arc=2pt, left=6pt, right=6pt, top=4pt, bottom=4pt,
  title={Hands-on Exercise}, fonttitle=\bfseries, coltitle=white,
  colbacktitle=orange!70!black
}
\newtcolorbox{homeworkbox}{
  colback=blue!6, colframe=blue!50!black, boxrule=0.6pt,
  arc=2pt, left=6pt, right=6pt, top=4pt, bottom=4pt,
  title={Homework}, fonttitle=\bfseries, coltitle=white,
  colbacktitle=blue!50!black
}

\title[\texttt{\textbackslash begin} to \texttt{\textbackslash end}]{From \texttt{\textbackslash begin} to \texttt{\textbackslash end}}
\subtitle{The Complete LaTeX Guide \\[2pt] \normalsize Day 1 of 7: Setup \& Document Anatomy}
\author[A.\ Rahman]{Atikur Rahman}
\institute[White Globe Initiative]{White Globe Initiative}

\setbeamertemplate{footline}{
  \leavevmode
  \hbox{%
  \begin{beamercolorbox}[wd=.38\paperwidth,ht=2.5ex,dp=1.125ex,leftskip=1em]{author in head/foot}%
    \usebeamerfont{author in head/foot}White Globe Initiative
  \end{beamercolorbox}%
  \begin{beamercolorbox}[wd=.42\paperwidth,ht=2.5ex,dp=1.125ex,center]{title in head/foot}%
    \usebeamerfont{title in head/foot}From \texttt{\textbackslash begin} to \texttt{\textbackslash end} --- Day 1
  \end{beamercolorbox}%
  \begin{beamercolorbox}[wd=.20\paperwidth,ht=2.5ex,dp=1.125ex,right,rightskip=1em]{date in head/foot}%
    \usebeamerfont{date in head/foot}\insertframenumber{} / \inserttotalframenumber
  \end{beamercolorbox}}%
  \vskip0pt
}

\begin{document}

\begin{frame}[plain]
  \begin{center}
    \vspace{4pt}
    {\Huge\bfseries\usebeamercolor[fg]{title}From \texttt{\textbackslash begin} to \texttt{\textbackslash end}}\\[6pt]
    {\Large\itshape The Complete LaTeX Guide}\\[18pt]
    {\large\bfseries Day 1 of 7 --- Setup \& Document Anatomy}\\[22pt]
    \rule{0.7\textwidth}{0.4pt}\\[14pt]
    \begin{tabular}{rl}
      \textbf{Instructor:}       & Atikur Rahman \\[3pt]
      \textbf{Organized by:}     & White Globe Initiative \\[3pt]
      \textbf{Instructor Email:} & \texttt{atik@whiteglobeinitiative.org} \\[3pt]
      \textbf{General Contact:}  & \texttt{contact@whiteglobeinitiative.org} \\[3pt]
      \textbf{Format:}           & 7-Day Intensive Course \\[3pt]
      \textbf{Session:}          & Day 1 of 7 \\
    \end{tabular}\\[18pt]
    \rule{0.7\textwidth}{0.4pt}
  \end{center}
\end{frame}

\begin{frame}{Today's Agenda}
  \tableofcontents
\end{frame}

\section{Why LaTeX?}

\begin{frame}{What Is LaTeX?}
  \small
  LaTeX is a \textbf{typesetting system}, not a word processor. Instead
  of clicking buttons to make text bold or centered, you write plain
  text with commands -- called \textbf{markup} -- that describe the
  \emph{structure} of your document, and a compiler turns that markup
  into a finished PDF. This is the opposite of Word or Google Docs,
  where what you see on screen \emph{is} the final output
  (\textbf{WYSIWYG}); in LaTeX, you describe what you \emph{mean}, and
  the system decides how it \emph{looks} (\textbf{WYSIWYM}).
  \vspace{6pt}
  \begin{center}
  \begin{tabular}{lcc}
    \toprule
    & \textbf{Word / Google Docs} & \textbf{LaTeX} \\
    \midrule
    Editing & Visual, direct & Source code + compile \\
    Equations & Clunky editor & Native, precise \\
    Tables \& refs & Manual & Automatic numbering \\
    Large documents & Fragile & Robust, modular \\
    \bottomrule
  \end{tabular}
  \end{center}
\end{frame}

\begin{frame}{Why Economists Use LaTeX}
  \small
  Four reasons economists adopt LaTeX over Word: \textbf{equations}
  (regressions, production functions, and game-theory payoffs are
  typeset natively and precisely, not pasted in as blurry images);
  \textbf{tables} (publication-quality regression tables in the exact
  style journals expect); \textbf{reproducibility} (plain-text source
  works cleanly with version control and co-author collaboration); and
  \textbf{journal requirements} (AER, JPE, QJE, and most Elsevier
  economics journals accept -- or require -- LaTeX and provide official
  templates, covered on Day 6). Bibliography management (BibTeX/natbib,
  Day 4) automates citations on top of all of this.
\end{frame}

\section{Getting Started with Overleaf}

\begin{frame}{Overleaf: Your LaTeX Environment}
  \small
  \textbf{Overleaf} (\texttt{overleaf.com}) is a free, browser-based
  LaTeX editor -- nothing to install, everything compiles in the cloud.
  The editor has two halves: a \textbf{source pane} on the left, where
  you type LaTeX code, and a \textbf{PDF preview} on the right, which
  updates every time you \textbf{recompile} (click the green button, or
  press Ctrl/Cmd+Enter). Today's workflow: \texttt{New Project}
  $\rightarrow$ \texttt{Blank Project} $\rightarrow$ type code
  $\rightarrow$ compile $\rightarrow$ read the result $\rightarrow$
  repeat. You will do this dozens of times today -- that loop \emph{is}
  how LaTeX is learned.
\end{frame}

\section{Anatomy of a Document}

\begin{frame}[fragile]{The Minimal LaTeX Document}
  \begin{lstlisting}
\documentclass{article}

\begin{document}
Hello, economics!
\end{document}
  \end{lstlisting}
  \begin{outputbox}
    \small Hello, economics!
  \end{outputbox}
  \whybox{Every LaTeX file has exactly this skeleton. The
  \textbf{preamble} (everything before \texttt{\textbackslash
  begin\{document\}}) configures the document -- here it just names
  the class. The \textbf{body} (between \texttt{\textbackslash
  begin\{document\}} and \texttt{\textbackslash end\{document\}}) is
  what actually gets typeset. Nothing outside \texttt{document} ever
  appears in your PDF.}
\end{frame}

\begin{frame}[fragile]{Document Class \& Packages}
  \begin{lstlisting}
\documentclass[11pt,a4paper]{article}
\usepackage[utf8]{inputenc}
\usepackage{amsmath}   % math tools (Day 2)
\usepackage{graphicx}  % images (Day 3)
  \end{lstlisting}
  \whybox{\texttt{\textbackslash documentclass\{article\}} tells LaTeX
  what \emph{kind} of document you're building -- \texttt{article} is
  correct for papers and working papers (\texttt{report} and
  \texttt{book} exist for longer documents). The options in
  \texttt{[...]} set font size and paper size. Each
  \texttt{\textbackslash usepackage\{...\}} loads one extra tool --
  almost everything useful in LaTeX (math, images, colors, citations)
  comes from a package, not the core language. This code has no
  visible output by itself -- it configures everything that follows.}
\end{frame}

\begin{frame}[fragile]{The Title Block}
  \begin{lstlisting}
\title{The Optimal Health-Investment Share}
\author{Atikur Rahman \\ White Globe Initiative}
\date{\today}

\begin{document}
\maketitle
  \end{lstlisting}
  \begin{outputbox}
    \begin{center}
      \normalsize\bfseries The Optimal Health-Investment Share\\[4pt]
      \small\normalfont Atikur Rahman\\ White Globe Initiative\\[2pt]
      \today
    \end{center}
  \end{outputbox}
  \whybox{\texttt{\textbackslash title}, \texttt{\textbackslash
  author}, and \texttt{\textbackslash date} are declared in the
  preamble -- they just \emph{store} the information. Nothing is drawn
  until \texttt{\textbackslash maketitle} is called inside the body,
  which is what actually renders the centered title block you see
  above. \texttt{\textbackslash today} auto-inserts the compile date.}
\end{frame}

\section{Structuring Your Paper}

\begin{frame}[fragile]{Sections and Subsections}
  \begin{lstlisting}
\section{Introduction}
\subsection{Motivation}

\section{Data and Methodology}
  \end{lstlisting}
  \begin{outputbox}
    \small
    \textbf{1\quad Introduction}\\[2pt]
    \hspace{10pt}\textbf{1.1\quad Motivation}\\[6pt]
    \textbf{2\quad Data and Methodology}
  \end{outputbox}
  \whybox{LaTeX \textbf{auto-numbers} every section and subsection in
  order -- notice how \texttt{Introduction} became ``1'' and
  \texttt{Motivation} became ``1.1'' with no numbers typed anywhere in
  the code. Insert or delete a section later and everything renumbers
  automatically. This structure also powers the table of contents and
  cross-referencing you'll use from Day 4 onward.}
\end{frame}

\begin{frame}[fragile]{Paragraphs and Line Breaks}
  \begin{lstlisting}
This is one paragraph. LaTeX ignores
single line breaks in your source code.

A blank line starts a new paragraph.
  \end{lstlisting}
  \begin{outputbox}
    \small This is one paragraph. LaTeX ignores single line breaks in
    your source code.\\[6pt]
    A blank line starts a new paragraph.
  \end{outputbox}
  \whybox{Notice the source has a line break after ``code'' but the
  output reads as one continuous sentence -- LaTeX treats a single
  line break as just a space. Only a truly \textbf{blank line} in your
  source starts a new paragraph (with the usual first-line indent). To
  force a line break \emph{without} starting a new paragraph, use
  \texttt{\textbackslash\textbackslash} instead.}
\end{frame}

\section{Text Formatting}

\begin{frame}[fragile]{Bold, Italics, and Comments}
  \begin{lstlisting}
This finding is \textbf{statistically significant}
at the \textit{5\% level}. % not the 10% level

Growth was \underline{negative} in 2020.
  \end{lstlisting}
  \begin{outputbox}
    \small This finding is \textbf{statistically significant} at the
    \textit{5\% level}.\\[4pt]
    Growth was \underline{negative} in 2020.
  \end{outputbox}
  \whybox{\texttt{\textbackslash textbf\{...\}} bolds, \texttt{
  \textbackslash textit\{...\}} italicizes, \texttt{\textbackslash
  underline\{...\}} underlines (rare in print economics papers -- use
  sparingly). The \texttt{\%} comment and everything after it on that
  line \textbf{never appears in the output} -- it's a note to yourself
  or a co-author. Special characters (\texttt{\% \$ \& \# \_}) need a
  backslash when you mean them literally, since LaTeX normally reads
  them as commands.}
\end{frame}

\section{Reading Compile Errors}

\begin{frame}[fragile]{Your First Compile Errors}
  \begin{lstlisting}
\textb{Significant results}
  \end{lstlisting}
  \begin{outputbox}
    \scriptsize\ttfamily
    ! Undefined control sequence.\\
    l.12 \textbackslash textb\{Significant results\}\\[4pt]
    \normalfont\normalsize No PDF is produced until this is fixed.
  \end{outputbox}
  \whybox{This is a typo: \texttt{\textbackslash textb} isn't a real
  command (it should be \texttt{\textbackslash textbf}). LaTeX stops
  and reports the exact line number -- in Overleaf this line is
  underlined in red, and clicking it jumps you straight there. The two
  other errors you'll meet constantly: \textit{Missing \$ inserted}
  (a math symbol used outside math mode -- fixed properly on Day 2) and
  a mismatched \texttt{\textbackslash begin}/\texttt{\textbackslash
  end} pair (an environment opened but never closed).}
\end{frame}

\section{Practice}

\begin{frame}{Hands-on Exercise}
  \begin{exercisebox}
    In a new Overleaf project, recreate the \textbf{title page and
    abstract} of a real, published economics paper:
    \begin{enumerate}
      \item Title, author, and date using \texttt{\textbackslash title},
        \texttt{\textbackslash author}, \texttt{\textbackslash date},
        \texttt{\textbackslash maketitle}.
      \item An \texttt{Abstract} section with one paragraph of body text.
      \item A first \texttt{\textbackslash section\{Introduction\}}
        heading.
      \item Bold at least one key term and italicize one phrase.
    \end{enumerate}
    Compile at least \textbf{five times}, intentionally breaking and
    fixing something each time (misspell a command, delete a brace,
    forget \texttt{\%}) -- and read what the resulting error says
    before you fix it.
  \end{exercisebox}
\end{frame}

\begin{frame}{Homework Before Day 2}
  \begin{homeworkbox}
    Pick the article you will rebuild all week (your own draft is
    ideal). Retype just the \textbf{introduction paragraph} into your
    Overleaf project -- no math or tables yet, those start tomorrow.
  \end{homeworkbox}
  \vspace{10pt}
  \small\textbf{Resources for today:}
  \begin{itemize}\small
    \item Overleaf -- ``Learn LaTeX in 30 Minutes''
    \item Overleaf -- ``Creating your first document in LaTeX''
    \item Overleaf -- ``Paragraphs and new lines'' / ``Bold, italics
      and underlining''
    \item Wikibooks LaTeX -- Getting Started chapter
  \end{itemize}
\end{frame}

\begin{frame}{Coming Up: Day 2}
  \Large
  \centering
  Math Mode for Economists \\[6pt]
  \normalsize
  Equations, subscripts, Greek letters, summations, and matrices --
  typesetting the models you actually work with.
\end{frame}

\begin{frame}[plain]
  \begin{center}
    \vspace{35pt}
    {\Large\bfseries Thank You}\\[10pt]
    {\large From \texttt{\textbackslash begin} to \texttt{\textbackslash end}: The Complete LaTeX Guide}\\[4pt]
    Day 1 of 7 --- Setup \& Document Anatomy\\[20pt]
    \rule{0.5\textwidth}{0.4pt}\\[10pt]
    Atikur Rahman \\
    White Globe Initiative \\[2pt]
    \texttt{atik@whiteglobeinitiative.org} \\
    \texttt{contact@whiteglobeinitiative.org}
  \end{center}
\end{frame}

\end{document}
